
Dynamics

A5) you only care about sinusoidal steady state dynamics, no transients

Control assumptions

Concept:
We introduce QLPS as method for optimal control of quasi-linearized systems subject to time-domain constraints.

We show that the QLPS for standard wave energy objectives and constraints is a QCQP.

We show that under additional conditions (narrow-band wave input and a singular PTO DOF), the QCQP reduces to order 1 and can be solved semi-ananlytically using geometry.

Assumptions:
QLPS requires assumptions A5 and A6:
\begin{itemize}
    \item A5)
    \textbf{Dynamics can be represented in linear form.}
    Assumption A2 quasi-linearizes the system, so this assumption is easily satisfied for MDOcean's dynamics.
    Rather than a linear state-space form, we choose a linear multiport form in which the state variables are ``efforts'' and ``flows'' which can be related with linear coefficient matrices because it is well suited for expressing power flow.

    \item A6)
    \textbf{All constraints can be written in terms of the fundamental of state variables}, in this case efforts and flows.
\end{itemize}
When additional assumptions A7 and A8 are satisfied, the QLPS problem can be represented as a QCQP:

\begin{itemize}
    \item A7)
    All constraints can be written with quadratic matrices.
    Examples of constraints that fit this restriction include effort amplitude, flow amplitude, average power, reactive power, and linear combinations thereof.

    \item A8)
    The objective function can be written as a quadratic metric.
    Average power is a quadratic metric.
\end{itemize}

When the additional assumptions A9 and A10 are satisfied, the QCQP can be solved semi-analytically:
\begin{itemize}
A9) \textbf{There is no coupling between dynamics at different frequencies.}
This is satisfied as a result of assumption A3, which states the wave excitation is narrow-band, and A2, which states that the WEC dynamics are low-pass.
A10) \textbf{There is no coupling between dynamics of different DOFs}, as viewed from the port where power is being maximized.
This is satisfied because the RM3 PTO contains only a single DOF.
Coupling between body DOFs is allowed, because they get collapsed into a single DOF at the PTO port.
\end{itemize}
Under these conditions, the QCQP for all frequencies and DOFs can be separated into a set of single-order QCQPs, one for each frequency and DOF, which can be solved semi-analytically using the geometric properties of circles.






Loadcase assumptions

A10) all nonzero JPD sea states are operations

A11) constructive interference with shaft torque

A12) use fundamental amplitude of fluid force (meh assumption bc not filtering)

A13) use fundamental amplitude for peak float and spar force (good bc filtering)

A14) surge force neglects diffraction